\chapter{Uruchomienie i badania prototypu}\label{ch:badania}
% Ok. 8-12 stron. Plan pomiarów: praca/plan-pracy.md sekcja 3 (9 pozycji).

\section{Stanowisko pomiarowe i metodyka}\label{sec:metodyka}
% - Sprzęt, warunki (J-Link odłączony, J3 pusty, DIN niski), niepewności.
\todo{Napisać.}

\section{Uruchamianie krok po kroku}\label{sec:uruchomienie}
% - Lutowanie U2, Z1 (BUCK2 bez nRF, R3 wylutowany), R1 SWD, pierwszy firmware.
\todo{Napisać.}

\section{Zasilanie i ładowanie}\label{sec:pomzasilanie}
% - Z1, Z2, Z4b, Ł1, Ł3, Ł5.
\todo{Napisać.}

\section{Radio}\label{sec:pomradio}
% - R1, R5 (RSSI w funkcji odległości, z płytą i bez).
\todo{Napisać.}

\section{Funkcje klawiatury}\label{sec:pomklawiatura}
% - K1-K2 na makiecie matrycy.
\todo{Napisać.}

\section{Pobór energii i czas pracy}\label{sec:pomenergia}
% - E1, E2.
\todo{Napisać.}

\section{Weryfikacja wymagań}\label{sec:weryfikacja}
% - Tabela W1..Wn -> test -> wynik -> spełnione/niespełnione.
\todo{Napisać.}

